\documentclass{scrartcl}
\usepackage{graphicx}  % required for inserting images

% for cyrillic and Bulgarian language
\usepackage[utf8]{inputenc}
\usepackage[T2A]{fontenc}
\usepackage[bulgarian]{babel}

% math packages
\usepackage{amsmath}
\usepackage{float}
\usepackage{amsfonts}
\usepackage{amssymb}
\usepackage{amsthm}
\usepackage{cancel}

\usepackage{ragged2e}  % adds FlushLeft

\usepackage{listings}
\usepackage{xcolor}

\definecolor{codegreen}{rgb}{0,0.6,0}
\definecolor{codegray}{rgb}{0.5,0.5,0.5}
\definecolor{codepurple}{rgb}{0.58,0,0.82}
\definecolor{backcolour}{rgb}{0.95,0.95,0.92}

\lstdefinestyle{mystyle}{
    label={lst:listing-cpp},
    language={C++},
    backgroundcolor=\color{backcolour},   
    commentstyle=\color{codegreen},
    keywordstyle=\color{magenta},
    numberstyle=\tiny\color{codegray},
    stringstyle=\color{codepurple},
    basicstyle=\ttfamily\footnotesize,
    breakatwhitespace=false,         
    breaklines=true,                 
    keepspaces=true,                  
    numbers=left,                    
    numbersep=5pt,                  
    showspaces=false,                
    showstringspaces=false,
    showtabs=false,                  
    tabsize=2,
}
\lstset{style=mystyle}

\usepackage[colorlinks=true, linkcolor=blue, urlcolor=cyan]{hyperref}

\newtheorem{theorem}{Теорема}
\newtheorem{definition}{Дефиниция}

\title{Компютърни мрежи и протоколи – OSI модел. Протоколи IPv4, IPv6, TCP, HTTP}
\author{Ивайло Андреев + gemini 3 Flash}
\date{31 май 2026}

\begin{document}
\maketitle

\tableofcontents
\newpage

\begin{FlushLeft}

\section{OSI модел – най-обща характеристика на нивата, съпоставяне с модела TCP/IP}

\subsection{OSI модел - най-обща характеристика на нивата}
\textbf{OSI (Open Systems Interconnection)} моделът е теоретичен абстрактен модел, описващ принципите на комуникация в мрежа и устройството на мрежовите протоколи. Той е въведен от организацията \textbf{ISO}. Основната градивна единица са слоевете (нивата), като всеки слой предоставя интерфейс и услуги към слоя над него. Моделът се състои от следните 7 нива, номерирани отдолу нагоре:

\begin{enumerate}
    \item \textbf{Физически слой (Physical Layer)}: Най-ниският слой на OSI модела. Отговаря за приемането и предаването на неструктурирани потоци от данни (битове) по физическия носител. Той осигурява физическата връзка между върховете в дадена топология. Неговите функции включват:
    \begin{itemize}
        \item \textit{Синхронизация на битовете} – осигурява механизми за синхронизация между изпращача и получателя.
        \item \textit{Честота на изпращане на битове (bit rate control)}.
        \item \textit{Определяне на мрежовата топология} (bus, star, mesh).
        \item \textit{Режим на предаване на битовете} – simplex, half-duplex или full-duplex.
        \item \textit{Начин на установяване и прекъсване на връзката между устройствата}.
        \item \textit{Кодиране и декодиране на информацията} (в двоичен вид) към и от електрически или физически сигнали.
    \end{itemize}
    Обектите на този слой са хардуерни устройства: мрежови карти (NIC), модеми, повторители (hubs) и коаксиални кабели.

    \item \textbf{Канален слой (Data Link Layer - DLL)}: Грижи се за безгрешното и надеждно предаване на данни от едно устройство до друго, утилизирайки физическия слой. Разделя се на два подслоя: \textbf{Logical Link Control (LLC)} и \textbf{Media Access Control (MAC)}. Той предава пакетите до правилните хостове, използвайки MAC адресите на устройствата, като за преобразуване между IP и MAC адреси често се използва \textbf{ARP (Address Resolution Protocol)}. Когато пакет от мрежовия слой пристигне, той бива разбиван на кадри (frames) спрямо размера, дефиниран от мрежовата карта. Функциите му са:
    \begin{itemize}
        \item \textit{Разбиване на кадри (framing)} – организира последователност чрез поставяне на специални битови последователности в началото и края на всеки кадър.
        \item \textit{Физическо адресиране} – добавяне на MAC адресите на изпращача/получателя в хедъра на кадъра.
        \item \textit{Контрол над грешките} – откриване и повторно предаване на счупени или загубени кадри.
        \item \textit{Контрол над потока (flow control)} – координация на количеството данни, предавани за единица време, поради различната скорост на устройствата.
        \item \textit{Контрол на достъпа} – MAC подслоят определя кой има контрол над споделения комуникационен канал в даден момент.
    \end{itemize}
    Имплементира се от Ethernet, суичове (switches) и бриджове (bridges).

    \item \textbf{Мрежови слой (Network Layer)}: Отговаря за предаването на данни между хостове от различни мрежи, които могат да се различават по своите физически и канални слоеве. Данните се капсулират под формата на пакети (\textbf{PDU – Protocol Data Unit}). Основни функции:
    \begin{itemize}
        \item \textit{Маршрутизация (Routing)} – намиране на най-кратък/оптимален път от изпращача до получателя през мрежата.
        \item \textit{Логическа адресация} – уникална идентификация на всяко устройство в глобалната мрежа чрез поставяне на IP адреси на изпращача и получателя в хедъра на пакета.
        \item \textit{Предотвратяване на натоварването (congestion control)} в мрежата.
    \end{itemize}
    Примерни протоколи: \textbf{IP (IPv4/IPv6)}, \textbf{ICMP}, \textbf{IPSec}.

    \item \textbf{Транспортен слой (Transport Layer)}: Осигурява доставка на съобщенията от край-до-край (end-to-end). Протоколите от този слой разбиват съобщенията, получени от сесийния слой, на по-малки единици, наречени \textbf{сегменти}, и ги подават на мрежовия слой. При получателя се грижи за реасемблирането им в правилната наредба и цялост. Функции:
    \begin{itemize}
        \item \textit{Сегментация и повторно сглобяване} – добавяне на метаданни в хедъра за реда на реасемблиране на сегментите.
        \item \textit{Адресация на услуги (Service Point Addressing)} – добавяне на порт на източника и дестинацията в хедъра на всеки сегмент, за да се достави съобщението до правилния процес.
    \end{itemize}
    Този слой се имплементира като част от операционната система. Примери: \textbf{TCP} (ориентиран към връзка, надежден) и \textbf{UDP} (без връзка, лек и бърз).

    \item \textbf{Сесиен слой (Session Layer)}: Отговаря за управлението на диалога между приложенията. Предоставя:
    \begin{itemize}
        \item \textit{Установяване, поддръжка и терминиране на сесии} между комуникиращи програми.
        \item \textit{Синхронизация} – добавяне на контролни точки (checkpoints) в потока от данни, което позволява възстановяване от мястото на прекъсване без загуба на информация.
        \item \textit{Диалогов контролер} – следи чий ред е да изпраща данни, поддържайки режими simplex, half-duplex или full-duplex.
        \item \textit{Сигурност} – извършване на процеси по автентикация.
    \end{itemize}
    Обичайно се имплементира чрез сокети (sockets).

    \item \textbf{Презентационен слой (Presentation Layer)}: Извлича и манипулира данните от приложния слой, грижейки се за синтаксиса и формата им, за да станат годни за пренос. Предоставя:
    \begin{itemize}
        \item \textit{Транслация} – конвертиране между различни формати за представяне на данни.
        \item \textit{Криптиране и декриптиране} – защита на информацията.
        \item \textit{Компресия} – намаляване на обема на пренасяните данни.
    \end{itemize}
    Примери: \textbf{TLS}, \textbf{SSL}, \textbf{MIME}.

    \item \textbf{Приложен слой (Application Layer)}: Най-горният слой, съвкупност от протоколи, които се имплементират директно от потребителските и сървърните процеси и приложения. Позволява на потребителя достъп до мрежови услуги. Примери: \textbf{HTTP}, \textbf{DNS}, \textbf{FTP}, \textbf{SMTP}.
\end{enumerate}

\subsection{Съпоставяне между моделите OSI и TCP/IP}
Докато OSI е теоретичен наръчник, описващ абстрактно функциите, моделът \textbf{TCP/IP} е разработен на базата на реално съществуващи стандартни протоколи. TCP/IP е по-компактна версия и се състои от 4 слоя:
\begin{enumerate}
    \item \textbf{Мрежов интерфейс (Network Interface / Достъп до мрежата)}: Обединява функциите на \textit{Физическия} и \textit{Каналния} слой на OSI.
    \item \textbf{Мрежов слой (Internet / Маршрутизация)}: Директно съответства на \textit{Мрежовия} слой на OSI.
    \item \textbf{Транспортен слой (Transport / Транспортиране)}: Съответства на \textit{Транспортния} слой на OSI.
    \item \textbf{Приложен слой (Application / Приложение)}: Обхваща и събира функциите на трите горни слоя на OSI – \textit{Сесиен}, \textit{Презентационен} и \textit{Приложен}.
\end{enumerate}

\textbf{Общи свойства}:
\begin{itemize}
    \item И двата модела се базират на единен стек от независими протоколи.
    \item Слоевете изпълняват сходни базови функции по обработка и транспорт на данните.
\end{itemize}

\textbf{Разлики}:
\begin{itemize}
    \item \textit{Разграничение на концепциите}: OSI прави ясно и строго разграничение между три основни свойства – услуги (какво прави слоят), интерфейси (как се достъпва) и протоколи (как се имплементира). При TCP/IP няма такова ясно разграничение.
    \item \textit{Време на създаване}: OSI моделът е измислен преди протоколите за него, което го прави по-абстрактен и универсален. Моделът TCP/IP е създаден след протоколите, нагоден е към тях и подмяната им е по-трудна.
    \item \textit{Връзка на транспортно ниво}: В класическия OSI модел мрежовото ниво поддържа connection-oriented режим, докато при TCP/IP интернет слоят е изцяло connectionless (IP), а транспортният слой предоставя избор между connectionless (UDP) и connection-oriented (TCP) режими.
\end{itemize}


\section{Протокол IPv4 адресация – класова и безкласова}

\textbf{IPv4} е фундаментален протокол от мрежовия слой, чиято основна задача е успешното предаване на пакети от източник до получател в рамките на глобалната мрежа. Това се постига чрез присвояване на уникален IP адрес на всеки хост и маршрутизатор.

\subsection{Структура на адреса}
Адресите в IPv4 са 32-битови двоични числа, които стандартно се записват като 4 десетични числа (октета), разделени с точка (напр. 192.0.2.1), като всеки октет е в интервала $[0, 255]$. Всеки адрес логически се състои от две части:
\begin{itemize}
    \item \textbf{Мрежов идентификатор (netid / префикс)} – идентифицира конкретната мрежа.
    \item \textbf{Идентификатор на хоста (hostid)} – идентифицира конкретното устройство в тази мрежа.
\end{itemize}

Специални адреси в хост частта:
\begin{itemize}
    \item Ако хост частта съдържа само двоични нули (\textbf{0}), това е \textit{адресът на самата мрежа}.
    \item Ако хост частта съдържа само единици (\textbf{1}), това е \textit{broadcast адресът} за едновременно предаване до всички хостове в тази мрежа.
    \item \textit{0.0.0.0} – обозначава неопределен адрес (unspecified address); записът \textit{0.0.0.0/0} се използва като маршрут по подразбиране (default route) в таблиците за маршрутизация.
\end{itemize}

\subsection{Класова адресация (Classful Addressing)}
В исторически план IPv4 адресите са били разделени на 5 фиксирани класа (A, B, C, D, E), характеризиращи се с твърд префикс и мрежова маска по подразбиране. 

Максималният брой използваеми хостове в дадена подмрежа се изчислява по следната математическа формула:
\begin{equation}
N_{\text{hosts}} = 2^{n} - 2
\end{equation}
където $n$ е броят на битовете, заделени за хостовата част на адреса.

\begin{itemize}
    \item \textbf{Клас A}:
    \begin{itemize}
        \item Фиксиран първи бит: \textbf{0}.
        \item Обхват на първия октет: \textbf{1.0.0.0} до \textbf{126.255.255.255} (Адрес 127.0.0.1 е заделен за loopback тест).
        \item Маска по подразбиране: \textbf{255.0.0.0} (/8 префикс).
        \item Капацитет: 127 мрежи с по над 16 милиона хоста всяка ($2^{24}-2$).
    \end{itemize}

    \item \textbf{Клас B}:
    \begin{itemize}
        \item Фиксирани първи битове: \textbf{10}.
        \item Обхват на първия октет: \textbf{128.0.0.0} до \textbf{191.255.255.255}.
        \item Маска по подразбиране: \textbf{255.255.0.0} (/16 префикс).
        \item Капацитет: около 16,000 мрежи с по 65,534 хоста всяка ($2^{16}-2$).
    \end{itemize}

    \item \textbf{Клас C}:
    \begin{itemize}
        \item Фиксирани първи битове: \textbf{110}.
        \item Обхват на първия октет: \textbf{192.0.0.0} до \textbf{223.255.255.255}.
        \item Маска по подразбиране: \textbf{255.255.255.0} (/24 префикс).
        \item Капацитет: над 2 милиона мрежи ($2^{21}$), но само с по 254 хоста всяка ($2^8-2$).
    \end{itemize}

    \item \textbf{Клас D}:
    \begin{itemize}
        \item Фиксирани първи битове: \textbf{1110}.
        \item Обхват: \textbf{224.0.0.0} до \textbf{239.255.255.255}.
        \item Предназначение: Резервиран изцяло за \textbf{multicasting} (групово предаване). Няма subnet маска и хост адреси, тъй като не адресира единичен хост.
    \end{itemize}

    \item \textbf{Клас E}:
    \begin{itemize}
        \item Фиксирани първи битове: \textbf{1111}.
        \item Обхват: \textbf{240.0.0.0} до \textbf{255.255.255.254}.
        \item Предназначение: Резервиран за експериментални цели и бъдещи изследвания. Няма subnet маска.
    \end{itemize}
\end{itemize}

\textbf{Недостатъци}: Класовата система се оказва изключително негъвкава. Половината от адресното пространство се заема от клас А за едва 127 системи. С разрастването на Интернет адресите от клас В се изчерпват светкавично, което принуждава раздаването на множество класове С на големи организации. Това води до драстично раздуване на глобалните маршрутизиращи (рутиращи) таблици, тъй като маршрутизаторите трябва да пазят записи за всяка отделна малка мрежа без възможност за агрегиране.

\subsection{Безкласова адресация (Classless Addressing - CIDR)}
За решаване на проблема с изчерпването на адресите и раздуването на таблиците се въвеждат техниките \textbf{VLSM (Variable-Length Subnet Masking)} и безкласовата адресация \textbf{CIDR (Classless Inter-Domain Routing)}.

При CIDR твърдите граници на класовете отпадат. Всеки адрес се представя като 32-битово число, последвано от наклонена черта „/“ и число, указващо точния брой на единиците в мрежовата маска (дължина на префикса). 
Например, в CIDR записа \textbf{192.0.2.96/28}, префиксът за мрежа (\textit{netid}) заема първите 28 бита, а за хостове остават $32 - 28 = 4$ бита ($2^4-2 = 14$ използваеми хоста). Съответната подмрежова маска е \textbf{255.255.255.240}. Това позволява мрежите да се разпределят с гъвкав размер и да се агрегират (супернетинг) с цел оптимизация на рутиращите таблици.


\section{Основни характеристики на протокол IPv6}

Поради окончателното изчерпване на IPv4 адресите, преходът към \textbf{IPv6} става наложителен, макар и бавен поради липсата на директна обратна съвместимост. IPv6 е проектиран да реши фундаменталните проблеми на своя предшественик.

\subsection{Ключови характеристики}
\begin{itemize}
    \item \textbf{Огромно адресно пространство}: Адресите в IPv6 са \textbf{128-битови}. Те предоставят практически неизчерпаем брой уникални адреси.
    \item \textbf{Опростен и разширяем хедър}: Много от тромавите и рядко използвани полета от IPv4 хедъра са премахнати или изнесени като опционални разширения (Extension Headers) в края на хедъра. Така базовият хедър остава с фиксиран и изчистен формат, което позволява двойно по-бърза обработка от рутерите, въпреки че самите адреси са 4 пъти по-дълги.
    \item \textbf{Комуникация от край-до-край (End-to-End)}: Всеки хост притежава глобално уникален публичен IP адрес. На практика това премахва необходимостта от NAT в повечето сценарии, възстановявайки оригиналния модел на директна комуникация, въпреки че в практиката все още съществуват форми на транслация като NAT66.
    \item \textbf{Автоматична конфигурация}: Поддържат се както \textit{stateful} (чрез DHCPv6), така и \textit{stateless} (\textbf{SLAAC – Stateless Address Autoconfiguration}) автоматична конфигурация на устройствата. Устройствата могат сами да генерират своя IP адрес без наличие на DHCP сървър.
    \item \textbf{Интегрирана сигурност}: Протоколът IPv6 е проектиран с вградена поддръжка на \textbf{IPSec}, осигуряващ рамка за автентикация и криптиране, въпреки че реалното му прилагане и използване не е задължително във всички съвременни софтуерни реализации.
    \item \textbf{Липса на Broadcast}: В IPv6 напълно отсъстват broadcast адресите. Тяхната функция се поема изцяло от по-ефективното \textbf{multicast} (групово) адресиране, предотвратявайки излишното натоварване на всички хостове в сегмента.
    \item \textbf{Anycast адресиране}: Един и същ IPv6 адрес може да бъде присвоен на група хостове (напр. разпръснати сървъри), като пакетите автоматично се рутират до най-близкия от тях.
    \item \textbf{Мобилност}: Позволява на хостовете физически и географски да сменят мрежата си, като на практика запазват текущия си IP адрес благодарение на автоконфигурацията и специфичните разширения в хедърите.
\end{itemize}

\subsection{Типове IPv6 адреси}
Адресите се записват като осем групи от по четири шестнадесетични цифри, разделени с двоеточие. За избягване на преливане извън текстовото поле при по-дълги адреси се използва безопасен режим на визуализация, например:

\verb|2001:0db8:0000:0000:0000:ff00:0042:8329|.

Основните дефинирани типове са:
\begin{itemize}
    \item \textbf{Неопределен адрес}: \verb|::/128| (всички битове са 0).
    \item \textbf{Loopback адрес}: \verb|::1/128| (еквивалент на 127.0.0.1).
    \item \textbf{Multicast}: \verb|FF00::/8| (започва с префикс \texttt{11111111} в двоичен вид).
    \item \textbf{Link-Local Unicast}: \verb|FE80::/10| – адреси за комуникация само в рамките на локалния физически сегмент.
    \item \textbf{Global Unicast}: Всички останали адреси – за стандартна глобална комуникация в Интернет.
\end{itemize}


\section{TCP – процедура на трикратно договаряне}

\textbf{TCP (Transmission Control Protocol)} е надежден, ориентиран към връзка (\textit{connection-oriented}) протокол от транспортния слой. Неговата цел е да осигури безгрешен двупосочен пренос на байтов поток от край-до-край, справяйки се успешно със загубени, дублирани или разбъркани пакети чрез механизма на поредни номера (\textit{sequence numbers}) в хедъра на сегментите. Той включва управление на потока (\textit{flow control}) чрез полето \textit{window size}.

Преди да започне какъвто и да е обмен на данни между клиента и сървъра, трябва задължително да се установи виртуална връзка чрез процедурата \textbf{трикратно договаряне (Three-Way Handshake)}. Тя преминава през три последователни стъпки:

\begin{enumerate}
    \item \textbf{SYN (Synchronize)}: Клиентът изпраща първоначален сегмент към сървъра, с който заявява желание за установяване на връзка. В този сегмент е вдигнат флагът \textbf{SYN}. Клиентът генерира и изпраща произволен начален пореден номер (\textit{Initial Sequence Number - ISN}), например $Seq = x$, с който ще номерира своите данни.
    
    \item \textbf{SYN + ACK (Synchronize-Acknowledge)}: Сървърът получава сегмента и ако е готов да приеме връзката, отговаря със сегмент, в който са вдигнати едновременно флаговете \textbf{SYN} и \textbf{ACK}. 
    \begin{itemize}
        \item Чрез флага \textbf{ACK} сървърът потвърждава получаването на клиентския SYN, като записва номер за потвърждение $Ack = x + 1$ (очаква следващия байт $x+1$).
        \item Чрез флага \textbf{SYN} сървърът указва своя собствен начален пореден номер за данните, които той ще предава в обратната посока, например $Seq = y$.
    \end{itemize}
    
    \item \textbf{ACK (Acknowledge)}: Клиентът получава отговора от сървъра и изпраща финален потвърдителен сегмент с вдигнат флаг \textbf{ACK}. Номерът за потвърждение се установява на $Ack = y + 1$ (с което се потвърждава получаването на SYN от сървъра), а поредният номер е $Seq = x + 1$.
\end{enumerate}

След успешното завършване на тази стъпка се установява \textit{full-duplex} връзка и двете страни могат да започнат реален и надежден обмен на приложни данни.


\section{Хипертекстов протокол HTTP}

\textbf{HTTP (Hypertext Transfer Protocol)} е приложен протокол, работещ на принципа клиент-сървър (заявка-отговор / \textit{request-response}). Обикновено клиентите са уеб браузъри, а сървърите са уеб сървъри, предоставящи достъп до обекти, наречени \textbf{ресурси}.

\subsection{Основни принципи и транспорт}
HTTP представлява чист текстов протокол. HTTP е \textbf{stateless} протокол (без състояние), което означава, че всяка заявка се разглежда сама за себе си и сървърът не пази информация за минали състояния и потребителски контекст между отделните заявки. За пренос на състояние на по-високо ниво се използват бисквитки (\textit{cookies}).

HTTP е концептуално независим от транспортния слой и може да работи върху всеки протокол, гарантиращ надежден транспорт на данни. В исторически план и при по-ранните си версии протоколът работи стандартно върху \textbf{TCP} сокети на порт \textbf{80} (или порт 443 за криптирания му вариант HTTPS). Съвременните версии на HTTP поддържат постоянни връзки (\textit{persistent connections}), което позволява един TCP сокет да се преизползва за няколко последователни заявки, но това е оптимизация и не нарушава stateless характера му.

В зависимост от конкретната версия обаче, технологичният стек за транспорт се променя:
\begin{itemize}
    \item \textbf{HTTP/1.1} използва базово \textbf{TCP}.
    \item \textbf{HTTP/2} също използва \textbf{TCP} с добавено мултиплексиране.
    \item Съвременният \textbf{HTTP/3} използва изцяло протокола \textbf{QUIC} върху \textbf{UDP} вместо TCP, с цел намаляване на закъснението при установяване на връзка и избягване на блокирането на линията (head-of-line blocking).
\end{itemize}

\subsection{Структура на съобщенията и методи}
Всяко съобщение (заявка или отговор) съдържа:
\begin{itemize}
    \item \textbf{Стартов ред} (указващ метода и Request URI при заявка, или статус код при отговор). Ресурсът се идентифицира чрез \textbf{URI} (Uniform Resource Identifier), който може да бъде определен по местоположение (\textbf{URL}) или по име (\textbf{URN}).
    \item \textbf{Хедъри (Headers)} – съдържат метаданни за съобщението.
    \item \textbf{Тяло (Response/Request Body)} – съдържа самото съдържание (HTML, изображения, JSON данни) и е отделено от хедърите с празен ред.
\end{itemize}

Протоколът дефинира методи за указване на семантиката на действието над ресурса. Основните методи са:
\begin{itemize}
    \item \textbf{GET}: Изисква репрезентация на конкретен ресурс. Трябва да се използва единствено и само за безопасно извличане на данни, без да променя състоянието им.
    \item \textbf{HEAD}: Напълно идентичен с \textbf{GET}, но изисква сървърът да върне само стартовия ред и хедърите (\textit{response header}), без да изпраща самото тяло (\textit{response body}). Използва се за бърза проверка на ресурса.
    \item \textbf{POST}: Изпраща данни (в \textit{request body}) към сървъра, които водят до промяна на съществуващ ресурс или създаване на нов подресурс.
    \item \textbf{PUT}: Заменя изцяло текущото представяне на указания ресурс с данните, изпратени в тялото на заявката (\textit{request body}).
    \item \textbf{PATCH}: Използва се за извършване на частична модификация над указания ресурс, базирана на изпратеното тяло.
    \item \textbf{DELETE}: Изтрива изрично указания в Request URI ресурс.
    \item \textbf{CONNECT}: Отваря мрежов тунел към сървъра, идентифициран от ресурса (често се използва при прокси сървъри).
    \item \textbf{OPTIONS}: Използва се за извличане и проверка на поддържаните HTTP методи от сървъра за конкретния URL.
    \item \textbf{TRACE}: Извършва loopback тест по пътя на заявката до ресурс, като кара сървъра да върне обратно точно това съобщение, което е получил, за диагностични цели.
\end{itemize}

\section*{Заключение}
Компютърните мрежи се изграждат върху йерархични архитектури от протоколи и услуги. OSI моделът предоставя теоретична рамка за разделяне на комуникационния процес на слоеве, докато TCP/IP представлява реално използваната архитектура в Интернет. Протоколите IPv4 и IPv6 осигуряват логическа адресация и маршрутизация, като IPv6 решава проблема с ограниченото адресно пространство на IPv4 чрез 128-битово адресиране и опростен хедър. TCP гарантира надежден транспорт чрез механизма на трикратното договаряне, а HTTP е основният приложен протокол за обмен на ресурси в Световната мрежа.

\end{FlushLeft}
\end{document}